
\chapter{2021年新高考II卷}

\section{单选题}

\begin{problem}
复数 \(\frac{2-\i}{1-3\i}\) 在复平面内对应的点所在的象限为（\quad）

\choices
    {第一象限}
    {第二象限}
    {第三象限}
    {第四象限}

\begin{answer}
A
\end{answer}
\end{problem}


\begin{problem}
若全集 \(U=\{1,2,3,4,5,6\}\)，\(A=\{1,3,6\}\)，\(B=\{2,3,4\}\)，则 \(A\cap \left( C_U B \right)=\)（\quad）

\choices
    {\(\{3\}\)}
    {\(\{1,6\}\)}
    {\(\{5,6\}\)}
    {\(\{1,3\}\)}

\begin{answer}
B
\end{answer}
\end{problem}


\begin{problem}
若抛物线 \(y^2=2px\left( p>0 \right)\) 的焦点到直线 \(y=x+1\) 的距离为 \(\sqrt2\)，则 \(p=\)（\quad）

\choices
    {\(1\)}
    {\(2\)}
    {\(2\sqrt2\)}
    {\(4\)}

\begin{answer}
B
\end{answer}
\end{problem}


\begin{problem}
北斗三号全球卫星导航系统是我国航天事业的重要成果. 在卫星导航系统中，地球静止同步卫星的轨道位于地球赤道所在平面，轨道高度为 \(36000\text{ km}\)（轨道高度是指卫星到地球表面的最短距离）. 把地球看成是一个球心为 \(O\)，半径 \(r\) 为 \(6400\text{ km}\) 的球，其上点 \(A\) 的纬度是指 \(OA\) 与赤道平面所成角的度数. 地球表面能直接观测到的一颗地球静止同步轨道卫星的点的纬度最大值记为 \(\alpha\)，该卫星信号覆盖地球表面的表面积为 \(S=2\pi r^2\left( 1-\cos \alpha \right)\)（单位：\(\text{km}^2\)），则 \(S\) 占地球表面积的百分比约为（\quad）

\choices
    {\(26\%\)}
    {\(34\%\)}
    {\(42\%\)}
    {\(50\%\)}

\begin{answer}
C
\end{answer}
\end{problem}


\begin{problem}
正四棱台的上、下底面的边长分别为 \(2\)，\(4\)，侧棱长为 \(2\)，则其体积为（\quad）

\choices
    {\(56\)}
    {\(28\sqrt2\)}
    {\(\frac{56}{3}\)}
    {\(\frac{28\sqrt2}{3}\)}

\begin{answer}
D
\end{answer}
\end{problem}


\begin{problem}
某物理量的测量结果服从正态分布 \(N\left( 10,\sigma^2 \right)\)，下列结论中不正确的是（\quad）

\choices
    {\(\sigma\) 越小，该物理量一次测量结果落在 \(\left( 9.9,10.1 \right)\) 内的概率越大}
    {该物理量一次测量结果大于 \(10\) 的概率为 \(0.5\)}
    {该物理量一次测量结果小于 \(9.99\) 与大于 \(10.01\) 的概率相等}
    {该物理量一次测量结果落在 \(\left( 9.9,10.2 \right)\) 与落在 \(\left( 10,10.3 \right)\) 内的概率相等}

\begin{answer}
D
\end{answer}
\end{problem}


\begin{problem}
若 \(a=\log_5 2\)，\(b=\log_8 3\)，\(c=\frac12\)，则下列判断正确的是（\quad）

\choices
    {\(c<b<a\)}
    {\(b<a<c\)}
    {\(a<c<b\)}
    {\(a<b<c\)}

\begin{answer}
C
\end{answer}
\end{problem}


\begin{problem}
已知函数 \(f\left( x \right)\) 的定义域为 \(\R\)，\(f\left( x+2 \right)\) 为偶函数，\(f\left( 2x+1 \right)\) 为奇函数，则（\quad）

\choices
    {\(f\left( -\frac12 \right)=0\)}
    {\(f\left( -1 \right)=0\)}
    {\(f\left( 2 \right)=0\)}
    {\(f\left( 4 \right)=0\)}

\begin{answer}
B
\end{answer}
\end{problem}

\section{多选题}

\begin{problem}
下列统计量中可用于度量样本 \(x_1\)，\(x_2\)，\(\cdots\)，\(x_n\) 离散程度的有（\quad）

\choices
    {样本 \(x_1\)，\(x_2\)，\(\cdots\)，\(x_n\) 的标准差}
    {样本 \(x_1\)，\(x_2\)，\(\cdots\)，\(x_n\) 的中位数}
    {样本 \(x_1\)，\(x_2\)，\(\cdots\)，\(x_n\) 的极差}
    {样本 \(x_1\)，\(x_2\)，\(\cdots\)，\(x_n\) 的平均数}

\begin{answer}
AC
\end{answer}
\end{problem}


\begin{problem}
如图，下列各正方体中，\(O\) 为下底面的中心，\(P\) 为所在棱的中点，\(M\)，\(N\) 为正方体的顶点. 则满足 \(MN\perp OP\) 的是（\quad）

\choices
    {\choicebitmap[width=0.15\paperwidth]{img/2021/new_gaokao_paper_2/q10_optA.png}}
    {\choicebitmap[width=0.15\paperwidth]{img/2021/new_gaokao_paper_2/q10_optB.png}}
    {\choicebitmap[width=0.15\paperwidth]{img/2021/new_gaokao_paper_2/q10_optC.png}}
    {\choicebitmap[width=0.15\paperwidth]{img/2021/new_gaokao_paper_2/q10_optD.png}}

\begin{answer}
BC
\end{answer}
\end{problem}


\begin{problem}
已知直线 \(l:ax+by-r^2=0\) 与圆 \(C:x^2+y^2=r^2\)，点 \(A\left( a,b \right)\)，则下列说法正确的是（\quad）

\choices
    {若点 \(A\) 在圆 \(C\) 上，则直线 \(l\) 与圆 \(C\) 相切}
    {若点 \(A\) 在圆 \(C\) 内，则直线 \(l\) 与圆 \(C\) 相离}
    {若点 \(A\) 在圆 \(C\) 外，则直线 \(l\) 与圆 \(C\) 相离}
    {若点 \(A\) 在直线 \(l\) 上，则直线 \(l\) 与圆 \(C\) 相切}

\begin{answer}
ABD
\end{answer}
\end{problem}


\begin{problem}
设正整数 \(n=a_0\cdot 2^0+a_1\cdot 2^1+\cdots +a_k\cdot 2^k\)，其中 \(a_i\in \{0,1\}\)，记 \(\omega\left( n \right)=a_0+a_1+\cdots +a_k\). 则（\quad）

\choices
    {\(\omega\left( 2n \right)=\omega\left( n \right)\)}
    {\(\omega\left( 2n+3 \right)=\omega\left( n \right)+1\)}
    {\(\omega\left( 8n+5 \right)=\omega\left( 4n+3 \right)\)}
    {\(\omega\left( 2^n-1 \right)=n\)}

\begin{answer}
ACD
\end{answer}
\end{problem}

\section{填空题}

\begin{problem}
已知双曲线 \(C:\frac{x^2}{a^2}-\frac{y^2}{b^2}=1\left( a>0,b>0 \right)\) 的离心率为 \(2\)，则 \(C\) 的渐近线方程为\fillinblank{}.

\begin{answer}
\(y=\pm\sqrt3x\)
\end{answer}
\end{problem}


\begin{problem}
写出一个函数 \(f\left( x \right)=\fillinblank{}\)，使其同时具有下列性质：

\begin{circlelist}
\item \(f\left( x_1x_2 \right)=f\left( x_1 \right)f\left( x_2 \right)\)；
\item 当 \(x\in \left( 0,+\infty \right)\) 时，\(f'\left( x \right)>0\)；
\item \(f'\left( x \right)\) 是奇函数.
\end{circlelist}

\begin{answer}
\(x^4\)
\end{answer}
\end{problem}


\begin{problem}
已知向量 \(\bs a+\bs b+\bs c=\bs 0\)，\(\left| \bs a \right|=1\)，\(\left| \bs b \right|=\left| \bs c \right|=2\)，则 \(\bs a\cdot \bs b+\bs b\cdot \bs c+\bs c\cdot \bs a=\fillinblank{}\).

\begin{answer}
\(-\frac{9}{2}\)
\end{answer}
\end{problem}


\begin{problem}
已知函数 \(f\left( x \right)=\left| \e^x-1 \right|\)，\(x_1<0\)，\(x_2>0\)，函数 \(f\left( x \right)\) 的图象在点 \(A\left( x_1,f\left( x_1 \right) \right)\) 和点 \(B\left( x_2,f\left( x_2 \right) \right)\) 处的两条切线互相垂直，且分别交 \(y\) 轴于 \(M\)，\(N\) 两点，则 \(\frac{\left| AM \right|}{\left| BN \right|}\) 取值范围是\fillinblank{}.

\begin{answer}
\((0,1)\)
\end{answer}
\end{problem}

\section{解答题}

\begin{problem}
记 \(S_n\) 是公差不为 \(0\) 的等差数列 \(\{a_n\}\) 的前 \(n\) 项和，若 \(a_3=S_5\)，\(a_2a_4=S_4\).

\begin{enumerate}
\item 求数列 \(\{a_n\}\) 的通项公式 \(a_n\)；
\item 求使 \(S_n>a_n\) 成立的 \(n\) 的最小值.
\end{enumerate}
\end{problem}

\begin{solution}
\begin{enumerate}
\item 设等差数列 \(\{a_n\}\) 的首项为 \(a_1=a\)，公差为 \(d\)，则 \(d\ne0\)，且 \(a_3=a+2d\). 由等差数列的前五项和等于五项中间项的五倍，得 \(S_5=5a_3\). 因为 \(a_3=S_5\)，所以 \(a_3=0\)，从而 \(a=-2d\). 因此 \(a_2=-d\)，\(a_4=d\)，且
\[
S_4=\frac{4}{2}\left(2a+3d\right)=2\left(-4d+3d\right)=-2d
\]
由 \(a_2a_4=S_4\)，得 \(-d^2=-2d\)，即 \(d(d-2)=0\). 由于 \(d\ne0\)，所以 \(d=2\)，\(a_1=-4\). 故数列的通项公式为 \(a_n=-4+2(n-1)=2n-6\).

\item 由通项公式得
\[
S_n=\frac{n\left(a_1+a_n\right)}{2}=\frac{n\left(-4+2n-6\right)}{2}=n(n-5)
\]
于是
\[
S_n>a_n\iff n(n-5)>2n-6\iff n^2-7n+6>0\iff (n-1)(n-6)>0
\]
因为 \(n\) 是正整数，所以满足条件的最小值为 \(n=7\).
\end{enumerate}
\end{solution}


\begin{problem}
在 \(\triangle ABC\) 中，角 \(A\)，\(B\)，\(C\) 所对的边长分别为 \(a\)，\(b\)，\(c\)，\(b=a+1\)，\(c=a+2\).

\begin{enumerate}
\item 若 \(2\sin C=3\sin A\)，求 \(\triangle ABC\) 的面积；
\item 是否存在正整数 \(a\)，使得 \(\triangle ABC\) 为钝角三角形？若存在，求出 \(a\) 的值；若不存在，说明理由.
\end{enumerate}
\end{problem}

\begin{solution}
\begin{enumerate}
\item 由正弦定理，\(\dfrac{\sin C}{\sin A}=\dfrac{c}{a}\). 结合已知条件 \(2\sin C=3\sin A\)，得
\[
\frac{c}{a}=\frac{3}{2}
\]
又因为 \(c=a+2\)，所以 \(\dfrac{a+2}{a}=\dfrac{3}{2}\)，解得 \(a=4\). 因而 \(b=5\)，\(c=6\). 设三角形的半周长为 \(s\)，则 \(s=\dfrac{4+5+6}{2}=\dfrac{15}{2}\). 由海伦公式，三角形面积为
\[
\sqrt{s(s-a)(s-b)(s-c)}=\sqrt{\frac{15}{2}\cdot\frac{7}{2}\cdot\frac{5}{2}\cdot\frac{3}{2}}=\frac{15\sqrt7}{4}
\]

\item 因为 \(c=a+2>b=a+1>a\)，所以 \(c\) 是最大边. 三角形存在的必要条件为 \(a+b>c\)，即 \(a+(a+1)>a+2\)，所以 \(a>1\). 当 \(a\) 为正整数时，\(a\ge2\). 三角形为钝角三角形当且仅当最大边满足
\[
c^2>a^2+b^2
\]
代入 \(b=a+1\)，\(c=a+2\)，得
\[
(a+2)^2>a^2+(a+1)^2\iff a^2-2a-3<0\iff (a-3)(a+1)<0
\]
结合 \(a>1\)，可得 \(1<a<3\). 因而正整数 \(a\) 只有 \(a=2\)，此时三边为 \(2\)，\(3\)，\(4\)，且 \(4^2>2^2+3^2\)，确为钝角三角形. 所以存在，且 \(a=2\).
\end{enumerate}
\end{solution}


\begin{problem}
在四棱锥 \(Q-ABCD\) 中，底面 \(ABCD\) 是正方形，若 \(AD=2\)，\(QD=QA=\sqrt5\)，\(QC=3\).

\begin{enumerate}
\item 证明：平面 \(QAD\perp \text{平面 }ABCD\)；
\item 求二面角 \(B-QD-A\) 的平面角的余弦值.
\end{enumerate}

\bitmapfigure[max width=0.34\linewidth]{img/2021/new_gaokao_paper_2/q19_fig1.png}
\end{problem}

\begin{solution}
\begin{enumerate}
\item 建立如图所示的空间直角坐标系，令底面 \(ABCD\) 位于 \(xOy\) 平面，取
\(A=(0,0,0)\)，\(B=(2,0,0)\)，\(C=(2,2,0)\)，\(D=(0,2,0)\).
设 \(Q=(x,y,z)\)，其中 \(z>0\). 由 \(QA=QD=\sqrt5\)，\(QC=3\)，得
\[
\begin{cases}
x^2+y^2+z^2=5,\\
x^2+(y-2)^2+z^2=5,\\
(x-2)^2+(y-2)^2+z^2=9.
\end{cases}
\]
前两式相减得 \(y=1\)，第三式与第一式相减得 \((x-2)^2-x^2=4\)，所以 \(x=0\). 再代入第一式得 \(z^2=4\)，故 \(z=2\). 设 \(H=(0,1,0)\)，则 \(H\) 在线段 \(AD\) 上，且 \(QH\) 垂直于底面 \(ABCD\). 又因为 \(QH\subset\) 平面 \(QAD\)，所以平面 \(QAD\perp\) 平面 \(ABCD\).

\item 在点 \(Q\) 处取向量 \(\overrightarrow{QD}=(0,1,-2)\)，\(\overrightarrow{QA}=(0,-1,-2)\)，\(\overrightarrow{QB}=(2,-1,-2)\).
二面角 \(B-QD-A\) 的平面角，是两个面内分别垂直于棱 \(QD\) 且分别指向 \(B\)，\(A\) 的向量的夹角. 将 \(\overrightarrow{QA}\)，\(\overrightarrow{QB}\) 分别在垂直于 \(QD\) 的方向上作正投影，注意到 \(|\overrightarrow{QD}|^2=5\)，\(\overrightarrow{QA}\cdot\overrightarrow{QD}=3\)，\(\overrightarrow{QB}\cdot\overrightarrow{QD}=3\).
所以
\[
\overrightarrow{QA_\perp}=\overrightarrow{QA}-\frac{\overrightarrow{QA}\cdot\overrightarrow{QD}}{|\overrightarrow{QD}|^2}\overrightarrow{QD}=\left(0,-\frac85,-\frac45\right)
\]
\[
\overrightarrow{QB_\perp}=\overrightarrow{QB}-\frac{\overrightarrow{QB}\cdot\overrightarrow{QD}}{|\overrightarrow{QD}|^2}\overrightarrow{QD}=\left(2,-\frac85,-\frac45\right)
\]
从而 \(\overrightarrow{QA_\perp}\cdot\overrightarrow{QB_\perp}=\frac{16}{5}\)，\(|\overrightarrow{QA_\perp}|=\frac{4}{\sqrt5}\)，\(|\overrightarrow{QB_\perp}|=\frac{6}{\sqrt5}\).
因此二面角的平面角的余弦值为
\[
\frac{\overrightarrow{QA_\perp}\cdot\overrightarrow{QB_\perp}}{|\overrightarrow{QA_\perp}|\,|\overrightarrow{QB_\perp}|}=\frac{16/5}{(4/\sqrt5)(6/\sqrt5)}=\frac23
\]
\end{enumerate}
\end{solution}

\begin{problem}
已知椭圆 \(C:\dfrac{x^2}{a^2}+\dfrac{y^2}{b^2}=1\left(a>b>0\right)\)，右焦点为 \(F\left(\sqrt2,0\right)\)，其离心率为 \(\dfrac{\sqrt6}{3}\).

\begin{enumerate}
\item 求椭圆 \(C\) 的方程；
\item 设 \(M\)，\(N\) 是椭圆 \(C\) 上的两点，直线 \(MN\) 与曲线 \(x^2+y^2=b^2\left(x>0\right)\) 相切. 证明：\(M\)，\(N\)，\(F\) 三点共线的充要条件是 \(|MN|=\sqrt3\).
\end{enumerate}
\end{problem}

\begin{solution}
\begin{enumerate}
\item 设椭圆的半焦距为 \(c\). 由右焦点坐标知 \(c=\sqrt2\)，由离心率公式 \(e=\dfrac ca\) 得
\[
a=\frac{\sqrt2}{\sqrt6/3}=\sqrt3
\]
又因为 \(b^2=a^2-c^2\)，所以 \(b^2=3-2=1\). 故椭圆 \(C\) 的方程为
\[
\frac{x^2}{3}+y^2=1
\]

\item 由上题知 \(b=1\). 因为直线 \(MN\) 与右半圆 \(x^2+y^2=1\left(x>0\right)\) 相切，所以可将其写成
直线 \(\ell:ux+vy=1\)，其中 \(u^2+v^2=1\)，\(u>0\).
其中切点为 \((u,v)\)，向量 \((-v,u)\) 是直线 \(\ell\) 的单位方向向量. 直线上的点可表示为
\[
(x,y)=(u,v)+t(-v,u)
\]
将其代入椭圆方程，得到关于 \(t\) 的方程
\[
\left(u^2+\frac{v^2}{3}\right)t^2+\frac{4uv}{3}t-\frac{2u^2}{3}=0
\]
设该方程的两个根为 \(t_1\)，\(t_2\)，它们对应交点 \(M\)，\(N\). 因为 \((-v,u)\) 是单位向量，所以
\[
|MN|=|t_1-t_2|
\]
该方程的判别式为
\[
\Delta=\left(\frac{4uv}{3}\right)^2-4\left(u^2+\frac{v^2}{3}\right)\left(-\frac{2u^2}{3}\right)=\frac{8u^2}{3}
\]
又因为 \(u^2+v^2=1\)，所以
\[
|MN|=\frac{\sqrt{\Delta}}{u^2+v^2/3}=\frac{2\sqrt6u}{1+2u^2}
\]

若 \(M\)，\(N\)，\(F\) 三点共线，则直线 \(\ell\) 经过 \(F\left(\sqrt2,0\right)\)，从而 \(u\sqrt2=1\)，即 \(u=\dfrac{1}{\sqrt2}\). 代入上式得 \(|MN|=\sqrt3\).

反之，若 \(|MN|=\sqrt3\)，则
\[
\frac{2\sqrt6u}{1+2u^2}=\sqrt3\iff 2\sqrt2u=1+2u^2\iff \left(\sqrt2u-1\right)^2=0
\]
所以 \(u=\dfrac{1}{\sqrt2}\)，于是 \(u\sqrt2=1\)，点 \(F\left(\sqrt2,0\right)\) 在直线 \(\ell\) 上，即 \(M\)，\(N\)，\(F\) 三点共线. 综上，\(M\)，\(N\)，\(F\) 三点共线的充要条件是 \(|MN|=\sqrt3\).
\end{enumerate}
\end{solution}

\begin{problem}
一种微生物群体可以经过自身繁殖不断生存下来，设一个这种微生物为第 \(0\) 代，经过一次繁殖后为第 \(1\) 代，再经过一次繁殖后为第 \(2\) 代，依此类推. 该微生物每代繁殖的个数是相互独立的且有相同的分布列，设 \(X\) 表示 \(1\) 个微生物个体繁殖下一代的个数，\(P(X=i)=p_i>0\)，其中 \(i=0\)，\(1\)，\(2\)，\(3\).

\begin{enumerate}
\item 已知 \(p_0=0.4\)，\(p_1=0.3\)，\(p_2=0.2\)，\(p_3=0.1\)，求 \(E(X)\)；
\item 设 \(p\) 表示该种微生物经过多代繁殖后临近灭绝的概率，\(p\) 是关于 \(x\) 的方程 \(p_0+p_1x+p_2x^2+p_3x^3=x\) 的一个最小正实根，求证：当 \(E(X)\le1\) 时，\(p=1\)，当 \(E(X)>1\) 时，\(p<1\)；
\item 根据你的理解说明（2）问结论的实际含义.
\end{enumerate}
\end{problem}

\begin{solution}
\begin{enumerate}
\item 由期望的定义，
\[
E(X)=0\cdot0.4+1\cdot0.3+2\cdot0.2+3\cdot0.1=1
\]

\item 令
\[
g(x)=p_0+p_1x+p_2x^2+p_3x^3-x
\]
则题中的方程就是 \(g(x)=0\)，并且 \(g(1)=0\). 对 \(x\ge0\)，有 \(g'(x)=p_1+2p_2x+3p_3x^2-1\)，\(g''(x)=2p_2+6p_3x>0\).
所以 \(g'(x)\) 在 \([0,+\infty)\) 上严格递增. 又因为 \(p_0+p_1+p_2+p_3=1\) 且 \(p_0\)，\(p_2\)，\(p_3>0\)，所以 \(g(0)=p_0>0\)，\(g'(0)=p_1-1<0\). 此外，\(g'(1)=p_1+2p_2+3p_3-1=E(X)-1\).

当 \(E(X)\le1\) 时，对任意 \(x\in[0,1)\)，由 \(g'\) 严格递增可得 \(g'(x)<g'(1)\le0\). 因而 \(g\) 在 \([0,1]\) 上递减，且对 \(x\in[0,1)\) 有 \(g(x)>g(1)=0\). 所以 \((0,1)\) 内没有方程的根，而 \(1\) 本身是正根，故最小正根 \(p=1\).

当 \(E(X)>1\) 时，\(g'(0)<0<g'(1)\). 由连续性及严格递增性，存在唯一的 \(\xi\in(0,1)\) 使 \(g'(\xi)=0\). 函数 \(g\) 在 \([0,\xi]\) 上递减，在 \([\xi,1]\) 上递增，因此 \(g(\xi)<g(1)=0\). 又因为 \(g(0)>0\)，由零点存在定理，存在 \(p_0'\in(0,\xi)\) 使 \(g(p_0')=0\). 于是方程有小于 \(1\) 的正根，故其最小正根满足 \(p<1\).

\item \(p\) 表示从一个微生物个体开始，经过许多代繁殖后群体最终临近灭绝的概率. 因而，当每个个体的平均繁殖数 \(E(X)\le1\) 时，群体最终临近灭绝的概率为 \(1\)，即几乎一定会灭绝；当 \(E(X)>1\) 时，临近灭绝的概率小于 \(1\)，所以群体有正概率能够持续繁殖而不灭绝.
\end{enumerate}
\end{solution}

\begin{problem}
已知函数 \(f(x)=(x-1)\e^x-ax^2+b\).

\begin{enumerate}
\item 讨论 \(f(x)\) 的单调性；
\item 从下面两个条件中任选一个作为已知条件，证明：\(f(x)\) 有一个零点.
\begin{circlelist}
\item \(\dfrac12<a\le\dfrac{\e^2}{2}\)，\(b>2a\)；
\item \(0<a<\dfrac12\)，\(b\le2a\).
\end{circlelist}
\end{enumerate}
\end{problem}

\begin{solution}
\begin{enumerate}
\item 求导得
\[
f'(x)=\e^x+(x-1)\e^x-2ax=x\left(\e^x-2a\right)
\]
当 \(a\le0\) 时，\(\e^x-2a>0\)，所以 \(f'(x)<0\)（\(x<0\)），\(f'(x)>0\)（\(x>0\)），函数在 \((-\infty,0)\) 上递减，在 \((0,+\infty)\) 上递增.

当 \(0<a<\dfrac12\) 时，\(\ln(2a)<0\). 因此
当 \(x<\ln(2a)\) 时 \(f'(x)>0\)；当 \(\ln(2a)<x<0\) 时 \(f'(x)<0\)；当 \(x>0\) 时 \(f'(x)>0\).
所以函数在 \((-\infty,\ln(2a))\) 上递增，在 \((\ln(2a),0)\) 上递减，在 \((0,+\infty)\) 上递增.

当 \(a=\dfrac12\) 时，\(f'(x)=x(\e^x-1)\ge0\)，且除 \(x=0\) 外为正，所以函数在 \(\R\) 上递增.

当 \(a>\dfrac12\) 时，\(\ln(2a)>0\). 因此
当 \(x<0\) 时 \(f'(x)>0\)；当 \(0<x<\ln(2a)\) 时 \(f'(x)<0\)；当 \(x>\ln(2a)\) 时 \(f'(x)>0\).
所以函数在 \((-\infty,0)\) 上递增，在 \((0,\ln(2a))\) 上递减，在 \((\ln(2a),+\infty)\) 上递增.

\item 选取第二个条件. 因为 \(0<a<\dfrac12\) 且 \(b\le2a\)，所以 \(b<1\)，从而
\[
f(0)=(-1)\e^0+b=b-1<0
\]
另一方面，\(a>0\)，且当 \(x\to+\infty\) 时，指数项 \((x-1)\e^x\) 的增长速度高于二次项 \(ax^2\)，故
\[
\lim_{x\to+\infty}f(x)=+\infty
\]
因此存在充分大的正数 \(T\)，使得 \(f(T)>0\). 由函数的连续性及零点存在定理，存在 \(x_0\in(0,T)\)，使得 \(f(x_0)=0\). 所以 \(f(x)\) 有一个零点.
\end{enumerate}
\end{solution}
